\documentclass[../DoAn.tex]{subfiles}
\begin{document}

\section{Tổng quan chương}
Chương 1 đã giới thiệu bối cảnh thực tiễn, hạn chế của các giải pháp hiện có và định hướng giải pháp của hệ thống Smart Grading. Dựa trên những phân tích đó, Chương 2 trình bày chi tiết kết quả khảo sát hiện trạng và phân tích yêu cầu của hệ thống. Nội dung chương gồm: tổng quan các sản phẩm tương tự và nhu cầu thực tế của người dùng; mô tả tổng quan chức năng của hệ thống thông qua biểu đồ use case tổng quát, biểu đồ phân rã cho từng nhóm chức năng chính và quy trình nghiệp vụ trọng yếu; đặc tả chi tiết cho các use case quan trọng; và cuối cùng là các yêu cầu phi chức năng làm cơ sở cho việc thiết kế và triển khai ở các chương sau.

\section{Khảo sát hiện trạng}
\label{section:2.1}

Để xác định rõ yêu cầu cho hệ thống Smart Grading, đồ án tiến hành khảo sát hiện trạng từ ba nguồn thông tin chính, bao gồm:
\begin{itemize}
    \item khảo sát nhu cầu của giáo viên và học sinh tại một số trường phổ thông và đại học;
    \item khảo sát hệ thống chấm thi đang được sử dụng tại trường;
    \item khảo sát các ứng dụng tương tự trên thị trường.
\end{itemize}

\noindent\textbf{Khảo sát nhu cầu người dùng.} Qua khảo sát trực tuyến và phỏng vấn sơ bộ khoảng 30 giáo viên và 80 học sinh, có thể nhận thấy các nhu cầu phổ biến như sau: giáo viên mong muốn có công cụ tạo đề thi nhanh từ ngân hàng câu hỏi có sẵn, hỗ trợ sinh nhiều mã đề để hạn chế gian lận, có khả năng chấm bài ngay sau khi kết thúc buổi thi để công bố kết quả sớm; học sinh mong muốn được xem điểm và bài làm của mình một cách minh bạch, có cơ chế gửi phúc khảo rõ ràng khi cho rằng điểm chưa chính xác. Ngoài ra, cả giáo viên và học sinh đều mong muốn hệ thống hỗ trợ thống kê kết quả theo lớp, theo môn, theo chủ đề để phục vụ cho việc đánh giá chất lượng học tập.

\noindent\textbf{Khảo sát hệ thống hiện tại.} Tại các trường được khảo sát, hệ thống chấm thi hiện đang sử dụng chủ yếu thuộc hai loại:
\begin{itemize}
    \item chấm thủ công bằng cách đối chiếu phiếu trả lời với đáp án, áp dụng phổ biến tại các trường có quy mô nhỏ;
    \item chấm bằng máy quét OMR chuyên dụng kết hợp phần mềm đi kèm, áp dụng tại một số trường lớn và các hội đồng thi cấp tỉnh.
\end{itemize}
Hệ thống loại thứ nhất có ưu điểm chi phí thấp nhưng tốn thời gian, dễ sai sót và khó thống kê. Hệ thống loại thứ hai cho tốc độ chấm nhanh và độ chính xác cao, nhưng chi phí đầu tư máy quét lớn, phần mềm thường đi kèm không hỗ trợ tốt các nghiệp vụ quản lý đào tạo và phúc khảo.

\noindent\textbf{Khảo sát ứng dụng tương tự.} Đồ án khảo sát một số giải pháp chấm thi phổ biến hiện nay, gồm: phần mềm thương mại như Remark Office OMR và Gradescope; công cụ mã nguồn mở như OMRChecker; các hệ thống thi trực tuyến như Google Forms, Microsoft Forms và Quizizz. Đặc điểm so sánh giữa các giải pháp được tổng hợp trong Bảng \ref{table:2.1}.

\begin{table}[H]
\centering
\begin{tabular}{|p{3cm}|p{3cm}|p{3cm}|p{3cm}|p{2cm}|}
\hline
\textbf{Tính năng} & \textbf{Remark OMR} & \textbf{Gradescope} & \textbf{OMRChecker} & \textbf{Smart Grading} \\ \hline
Chấm OMR tự động & Có & Có & Có & Có \\ \hline
Tạo nhiều mã đề & Không & Không & Không & Có \\ \hline
Quản lý lớp, học sinh & Không & Có & Không & Có \\ \hline
Phúc khảo trực tuyến & Không & Có & Không & Có \\ \hline
Chấm trên di động & Không & Có & Có & Có \\ \hline
Hỗ trợ tiếng Việt & Không & Không & Một phần & Có \\ \hline
Tích hợp AI & Không & Có & Không & Có \\ \hline
Chi phí đầu tư & Cao & Cao & Miễn phí & Miễn phí \\ \hline
\end{tabular}
\caption{So sánh tính năng giữa các giải pháp chấm thi hiện có}
\label{table:2.1}
\end{table}

Qua bảng so sánh có thể nhận thấy hệ thống Smart Grading được đề xuất trong đồ án có sự kết hợp đầy đủ các tính năng mà các giải pháp hiện có chưa đáp ứng đồng thời, đặc biệt là khả năng tạo nhiều mã đề, quản lý đào tạo tích hợp, hỗ trợ phúc khảo minh bạch, chấm trên di động và hỗ trợ tiếng Việt đầy đủ.

\section{Tổng quan chức năng}
\label{section:2.2}

Trên cơ sở kết quả khảo sát, hệ thống Smart Grading được đề xuất với các nhóm chức năng chính gồm: quản lý người dùng và đào tạo (trường, lớp, giáo viên, học sinh), quản lý ngân hàng câu hỏi và đề thi, chấm thi bằng nhận dạng OMR, quản lý bài thi và phúc khảo, và báo cáo thống kê. Biểu đồ use case tổng quát mô tả các tác nhân và use case chính của hệ thống như trong Hình \ref{fig:usecase_tongquat}.

\begin{figure}[H]
    \centering
    \includegraphics[width=0.85\textwidth]{Hinhve/UseCaseTongQuat.png}
    \caption{Biểu đồ use case tổng quát của hệ thống Smart Grading}
    \label{fig:usecase_tongquat}
\end{figure}

\noindent\textbf{Tác nhân tham gia hệ thống.} Hệ thống có bốn tác nhân chính:
\begin{itemize}
    \item \textit{Quản trị viên hệ thống} (admin): chịu trách nhiệm quản lý toàn bộ trường, tài khoản và cấu hình hệ thống;
    \item \textit{Quản trị viên trường} (school-admin): quản lý các lớp, giáo viên, học sinh thuộc trường mình phụ trách;
    \item \textit{Giáo viên} (teacher): thực hiện tạo đề thi, quản lý ngân hàng câu hỏi, chấm bài và xử lý phúc khảo;
    \item \textit{Học sinh} (student): xem đề thi, xem điểm, gửi yêu cầu phúc khảo và sử dụng trợ lý AI.
\end{itemize}

\noindent\textbf{Các use case chính.} Hệ thống có 13 use case chính, được phân thành 4 nhóm chức năng lớn và sẽ được phân rã chi tiết ở các mục tiếp theo:
\begin{itemize}
    \item \textit{Nhóm quản lý đào tạo:} Quản lý trường học, Quản lý giáo viên, Quản lý học sinh, Đăng ký/Đăng nhập;
    \item \textit{Nhóm ngân hàng câu hỏi và đề thi:} Tạo câu hỏi/đề thi (có tương tác System), Tạo đề thi và trộn câu cho nhiều đề thi;
    \item \textit{Nhóm chấm thi và phúc khảo:} Quản lý các bài thi của học sinh, Chấm thi bằng ứng dụng, Tạo đơn phúc khảo, Phản hồi phúc khảo, Xem bài thi và điểm số;
    \item \textit{Nhóm thống kê:} Tạo báo cáo cho bài kiểm tra (có tương tác System).
\end{itemize}

\subsection{Biểu đồ use case phân rã Đăng ký/Đăng nhập}
\label{subsection:2.2.1}

Use case Đăng ký/Đăng nhập bao gồm các use case con:
\begin{itemize}
    \item Đăng ký tài khoản;
    \item Đăng nhập hệ thống;
    \item Xác thực email;
    \item Đặt lại mật khẩu;
    \item Đăng xuất;
    \item Quản lý phiên truy cập.
\end{itemize}
Tác nhân chính là tất cả bốn vai trò người dùng.

\subsection{Biểu đồ use case phân rã Quản lý trường học}
\label{subsection:2.2.2}

Use case Quản lý trường học gồm:
\begin{itemize}
    \item Tạo trường mới;
    \item Cập nhật thông tin trường;
    \item Xóa trường;
    \item Xem danh sách trường;
    \item Cấu hình thang điểm và tham số chấm.
\end{itemize}
Tác nhân chính là Quản trị viên hệ thống.

\subsection{Biểu đồ use case phân rã Quản lý giáo viên}
\label{subsection:2.2.3}

Use case Quản lý giáo viên gồm:
\begin{itemize}
    \item Thêm giáo viên;
    \item Cập nhật thông tin giáo viên;
    \item Phân công giáo viên theo môn;
    \item Vô hiệu hóa tài khoản giáo viên;
    \item Xem danh sách giáo viên.
\end{itemize}
Tác nhân chính là Quản trị viên hệ thống và Quản trị viên trường.

\subsection{Biểu đồ use case phân rã Quản lý học sinh}
\label{subsection:2.2.4}

Use case Quản lý học sinh gồm:
\begin{itemize}
    \item Thêm học sinh vào lớp;
    \item Cập nhật thông tin học sinh;
    \item Chuyển lớp cho học sinh;
    \item Xóa học sinh khỏi lớp;
    \item Xem danh sách học sinh theo lớp.
\end{itemize}
Tác nhân chính là Quản trị viên trường và Giáo viên.

\subsection{Biểu đồ use case phân rã Tạo câu hỏi, đề thi}
\label{subsection:2.2.5}

Use case Tạo câu hỏi và đề thi có sự tương tác của System, gồm:
\begin{itemize}
    \item Tạo câu hỏi thủ công;
    \item Sinh câu hỏi bằng AI;
    \item Phê duyệt câu hỏi;
    \item Nhập khẩu câu hỏi từ file;
    \item Tạo bộ đề thi;
    \item Cấu hình tham số đề thi.
\end{itemize}
Tác nhân chính là Giáo viên, có sự hỗ trợ của System thông qua các dịch vụ AI.

\subsection{Biểu đồ use case phân rã Tạo đề thi, trộn các câu cho nhiều đề thi}
\label{subsection:2.2.6}

Use case Tạo đề thi và trộn câu cho nhiều đề thi gồm:
\begin{itemize}
    \item Chọn câu hỏi từ ngân hàng;
    \item Cấu hình số lượng phiên bản;
    \item Thiết lập cấu hình trộn (câu hỏi, đáp án);
    \item Sinh phiếu OMR cho mỗi phiên bản;
    \item Phát hành đề thi;
    \item Xuất bản đề thi sang LaTeX/PDF.
\end{itemize}
Tác nhân chính là Giáo viên.

\subsection{Biểu đồ use case phân rã Quản lý các bài thi của học sinh}
\label{subsection:2.2.7}

Use case Quản lý các bài thi của học sinh gồm:
\begin{itemize}
    \item Ghi nhận bài nộp từ thiết bị di động;
    \item Xem danh sách bài nộp;
    \item Chỉnh sửa thủ công kết quả chấm;
    \item Phê duyệt kết quả bài làm;
    \item Xóa bài nộp.
\end{itemize}
Tác nhân chính là Giáo viên.

\subsection{Biểu đồ use case phân rã Chấm thi bằng ứng dụng}
\label{subsection:2.2.8}

Use case Chấm thi bằng ứng dụng gồm:
\begin{itemize}
    \item Chụp ảnh phiếu trả lời bằng camera;
    \item Tiền xử lý ảnh (làm phẳng, cắt, nhị phân hóa);
    \item Nhận dạng mã đề và mã học sinh;
    \item Nhận dạng các lựa chọn được tô;
    \item Đối chiếu đáp án và tính điểm;
    \item Đồng bộ kết quả lên server.
\end{itemize}
Tác nhân chính là Giáo viên thông qua thiết bị di động.

\subsection{Biểu đồ use case phân rã Tạo đơn phúc khảo}
\label{subsection:2.2.9}

Use case Tạo đơn phúc khảo gồm:
\begin{itemize}
    \item Chọn bài thi cần phúc khảo;
    \item Chọn câu hỏi cần phúc khảo;
    \item Nhập lý do phúc khảo;
    \item Đính kèm hình ảnh bằng chứng (nếu có);
    \item Gửi đơn phúc khảo;
    \item Theo dõi trạng thái đơn.
\end{itemize}
Tác nhân chính là Học sinh.

\subsection{Biểu đồ use case phân rã Phản hồi phúc khảo}
\label{subsection:2.2.10}

Use case Phản hồi phúc khảo gồm:
\begin{itemize}
    \item Tiếp nhận đơn phúc khảo;
    \item Xem chi tiết bài làm của học sinh;
    \item So sánh với đáp án;
    \item Phê duyệt hoặc từ chối;
    \item Nhập phản hồi cho học sinh;
    \item Cập nhật lại điểm (nếu phê duyệt).
\end{itemize}
Tác nhân chính là Giáo viên.

\subsection{Biểu đồ use case phân rã Xem bài thi và điểm số}
\label{subsection:2.2.11}

Use case Xem bài thi và điểm số gồm:
\begin{itemize}
    \item Xem danh sách bài thi;
    \item Xem điểm tổng và điểm từng câu;
    \item Xem ảnh phiếu đã chấm;
    \item Xem đáp án đúng;
    \item Tải về báo cáo điểm cá nhân.
\end{itemize}
Tác nhân chính là Học sinh.

\subsection{Biểu đồ use case phân rã Tạo báo cáo cho bài kiểm tra}
\label{subsection:2.2.12}

Use case Tạo báo cáo cho bài kiểm tra có sự tương tác của System, gồm:
\begin{itemize}
    \item Chọn bài kiểm tra cần báo cáo;
    \item Yêu cầu hệ thống tổng hợp dữ liệu;
    \item System phân tích thống kê (phổ điểm, độ khó câu hỏi, độ phân biệt);
    \item Sinh báo cáo dạng biểu đồ và bảng số liệu;
    \item Xuất báo cáo ra PDF/Excel.
\end{itemize}
Tác nhân chính là Giáo viên, có sự hỗ trợ của System thông qua các dịch vụ phân tích.

\subsection{Quy trình nghiệp vụ}
\label{subsection:2.2.13}

Hệ thống Smart Grading có một số quy trình nghiệp vụ quan trọng kết hợp nhiều use case với nhau. Quy trình nghiệp vụ "Tổ chức kỳ thi và chấm bài" được mô tả tổng quan như sau: đầu tiên giáo viên tạo đề thi từ ngân hàng câu hỏi, hệ thống tự động trộn đề thành nhiều phiên bản và sinh phiếu OMR tương ứng; giáo viên in phiếu và phát cho học sinh làm bài; sau khi thu phiếu, giáo viên sử dụng ứng dụng di động để chấm từng phiếu; kết quả chấm được đồng bộ về server và công bố cho học sinh; học sinh có thể xem điểm và gửi đơn phúc khảo nếu cần; giáo viên tiếp nhận và phản hồi phúc khảo; cuối cùng giáo viên tạo báo cáo thống kê cho bài kiểm tra.

Quy trình nghiệp vụ "Quản lý đào tạo" gồm: admin tạo trường học và tài khoản quản trị trường; quản trị trường tạo năm học, khối lớp, môn học và tài khoản giáo viên; giáo viên được phân công phụ trách lớp và môn học, tạo lớp và thêm học sinh vào lớp; học sinh được cấp tài khoản và đăng nhập vào hệ thống để xem thông tin cá nhân, kết quả học tập.

Quy trình nghiệp vụ "Phúc khảo" gồm: học sinh xem bài thi và phát hiện điểm chưa chính xác, gửi đơn phúc khảo kèm lý do và bằng chứng; giáo viên nhận thông báo, mở chi tiết bài làm của học sinh, so sánh với đáp án; giáo viên phê duyệt hoặc từ chối đơn phúc khảo kèm phản hồi; nếu phê duyệt, điểm của học sinh được cập nhật lại và thông báo kết quả được gửi cho học sinh.

\section{Đặc tả chức năng}
\label{section:2.3}

Căn cứ vào tầm quan trọng và độ phức tạp, đồ án lựa chọn đặc tả chi tiết cho 6 use case chính gồm:
\begin{itemize}
    \item Đăng ký/Đăng nhập;
    \item Tạo đề thi, trộn câu cho nhiều đề thi;
    \item Chấm thi bằng ứng dụng;
    \item Tạo đơn phúc khảo;
    \item Phản hồi phúc khảo;
    \item Tạo báo cáo cho bài kiểm tra.
\end{itemize}
Mỗi đặc tả gồm tên use case, luồng sự kiện chính, luồng sự kiện phát sinh, tiền điều kiện và hậu điều kiện.

\subsection{Đặc tả use case Đăng ký/Đăng nhập}
\label{subsection:2.3.1}

\noindent\textbf{Tên use case:} Đăng ký/Đăng nhập.

\noindent\textbf{Tiền điều kiện:} Người dùng truy cập được vào ứng dụng web hoặc ứng dụng di động Smart Grading.

\noindent\textbf{Luồng sự kiện chính:}
\begin{enumerate}
    \item Người dùng truy cập trang đăng nhập và chọn chức năng "Đăng ký" nếu chưa có tài khoản.
    \item Hệ thống hiển thị form đăng ký yêu cầu nhập họ tên, email, mật khẩu và xác nhận mật khẩu.
    \item Người dùng nhập thông tin và gửi yêu cầu đăng ký.
    \item Hệ thống kiểm tra tính hợp lệ của email, độ mạnh của mật khẩu và sự trùng lặp email trong cơ sở dữ liệu.
    \item Hệ thống tạo tài khoản mới với trạng thái chưa xác thực email, gửi email xác thực chứa mã token đến email của người dùng.
    \item Người dùng truy cập email, nhấp vào liên kết xác thực.
    \item Hệ thống đánh dấu tài khoản đã xác thực email và hiển thị thông báo xác thực thành công.
    \item Người dùng quay lại trang đăng nhập, nhập email và mật khẩu để đăng nhập.
    \item Hệ thống kiểm tra thông tin đăng nhập, cấp access token và refresh token cho người dùng.
    \item Hệ thống chuyển hướng người dùng đến trang dashboard tương ứng với vai trò của họ.
\end{enumerate}

\noindent\textbf{Luồng sự kiện phát sinh:}
\begin{itemize}
    \item Tại bước 4, nếu email đã tồn tại: hệ thống thông báo email đã được đăng ký và đề xuất chức năng đăng nhập hoặc đặt lại mật khẩu.
    \item Tại bước 4, nếu mật khẩu không đủ mạnh: hệ thống yêu cầu người dùng nhập lại mật khẩu theo yêu cầu bảo mật.
    \item Tại bước 7, nếu liên kết xác thực đã hết hạn: hệ thống cho phép gửi lại email xác thực mới.
    \item Tại bước 9, nếu thông tin đăng nhập sai: hệ thống thông báo sai thông tin và cho phép thử lại tối đa 5 lần trước khi khóa tài khoản tạm thời.
\end{itemize}

\noindent\textbf{Hậu điều kiện:} Nếu thành công, người dùng được đăng nhập vào hệ thống và có quyền truy cập các chức năng tương ứng với vai trò. Nếu thất bại, trạng thái hệ thống không thay đổi.

\subsection{Đặc tả use case Tạo đề thi, trộn các câu cho nhiều đề thi}
\label{subsection:2.3.2}

\noindent\textbf{Tên use case:} Tạo đề thi, trộn các câu cho nhiều đề thi.

\noindent\textbf{Tiền điều kiện:} Giáo viên đã đăng nhập vào hệ thống, đã có ngân hàng câu hỏi thuộc môn học tương ứng.

\noindent\textbf{Luồng sự kiện chính:}
\begin{enumerate}
    \item Giáo viên chọn chức năng "Tạo đề thi" trên giao diện web.
    \item Hệ thống hiển thị form tạo đề thi với các trường: tiêu đề, mô tả, lớp áp dụng, môn học, thời gian làm bài, tổng điểm, số lượng phiên bản cần tạo.
    \item Giáo viên nhập thông tin chung và lưu bản nháp.
    \item Giáo viên chọn câu hỏi từ ngân hàng câu hỏi theo chủ đề, độ khó hoặc theo danh sách đã chuẩn bị trước.
    \item Giáo viên cấu hình tham số trộn đề: có trộn thứ tự câu hỏi hay không, có trộn thứ tự đáp án hay không, số lượng phiên bản cần tạo (tối đa 50).
    \item Giáo viên nhấn nút "Sinh đề thi" để yêu cầu hệ thống tạo các phiên bản đề.
    \item Hệ thống tạo các phiên bản đề thi bằng cách hoán vị câu hỏi và đáp án theo cấu hình, đảm bảo đáp án đúng được cập nhật tương ứng cho từng phiên bản.
    \item Hệ thống sinh mã nguồn LaTeX sử dụng gói AMC cho từng phiên bản đề, biên dịch sang PDF và tải lên dịch vụ lưu trữ đám mây.
    \item Hệ thống sinh phiếu trả lời OMR kèm tọa độ nhận dạng cho từng phiên bản, lưu thông tin tọa độ vào cơ sở dữ liệu.
    \item Giáo viên xem lại danh sách các phiên bản đề đã tạo, tải về file PDF đề thi và phiếu trả lời.
    \item Giáo viên nhấn nút "Phát hành" để công bố đề thi cho học sinh.
    \item Hệ thống gửi thông báo đến học sinh thuộc các lớp áp dụng và cập nhật trạng thái đề thi thành "đã phát hành".
\end{enumerate}

\noindent\textbf{Luồng sự kiện phát sinh:}
\begin{itemize}
    \item Tại bước 3, nếu thiếu thông tin bắt buộc: hệ thống yêu cầu giáo viên bổ sung.
    \item Tại bước 4, nếu ngân hàng câu hỏi không đủ số lượng cần thiết: hệ thống cảnh báo và đề xuất giảm số lượng phiên bản.
    \item Tại bước 7, nếu có lỗi khi trộn đề: hệ thống thông báo và giữ nguyên bản nháp.
    \item Tại bước 8, nếu biên dịch LaTeX thất bại: hệ thống ghi nhận log lỗi và cho phép giáo viên thử lại.
\end{itemize}

\noindent\textbf{Hậu điều kiện:} Nếu thành công, đề thi được tạo với nhiều phiên bản, phiếu OMR tương ứng được sinh và học sinh nhận được thông báo. Nếu thất bại, trạng thái đề thi giữ nguyên ở bản nháp.

\subsection{Đặc tả use case Chấm thi bằng ứng dụng}
\label{subsection:2.3.3}

\noindent\textbf{Tên use case:} Chấm thi bằng ứng dụng.

\noindent\textbf{Tiền điều kiện:} Giáo viên đã đăng nhập trên ứng dụng di động, có bài thi đã được phát hành và phiếu OMR tương ứng.

\noindent\textbf{Luồng sự kiện chính:}
\begin{enumerate}
    \item Giáo viên mở ứng dụng di động Smart Grading và chọn chức năng "Quét phiếu".
    \item Ứng dụng yêu cầu quyền truy cập camera và hiển thị giao diện chụp ảnh.
    \item Giáo viên đặt phiếu trả lời trên bề mặt phẳng, căn chỉnh vào khung hình và chụp ảnh.
    \item Ứng dụng tải thông tin mẫu OMR tương ứng với bài thi từ server.
    \item Ứng dụng thực hiện tiền xử lý ảnh: chuyển ảnh sang thang xám, làm mờ Gaussian, nhị phân hóa bằng phương pháp Otsu, phát hiện và chỉnh nghiêng phiếu, cắt vùng phiếu.
    \item Ứng dụng phát hiện các ô trả lời bằng cách so sánh vị trí với tọa độ đã lưu trong mẫu OMR.
    \item Ứng dụng đánh giá mức độ tô của từng ô, xác định ô được chọn dựa trên ngưỡng diện tích tô, phát hiện lỗi tô nhầm (nhiều ô trong cùng một câu) hoặc tô nhẹ.
    \item Ứng dụng nhận dạng mã đề và mã số sinh viên từ vùng đầu phiếu.
    \item Ứng dụng đối chiếu câu trả lời của học sinh với đáp án của phiên bản đề tương ứng, tính điểm cho từng câu và tổng điểm.
    \item Ứng dụng hiển thị kết quả chấm cho giáo viên kèm ảnh phiếu đã được đánh dấu.
    \item Giáo viên kiểm tra và có thể chỉnh sửa thủ công các câu bị nhận dạng sai hoặc đánh dấu là không hợp lệ.
    \item Giáo viên nhấn nút "Lưu kết quả" để đồng bộ kết quả lên server.
\end{enumerate}

\noindent\textbf{Luồng sự kiện phát sinh:}
\begin{itemize}
    \item Tại bước 5, nếu ảnh chụp mờ hoặc nghiêng quá nhiều: ứng dụng yêu cầu giáo viên chụp lại.
    \item Tại bước 8, nếu không nhận dạng được mã đề: ứng dụng cho phép giáo viên chọn thủ công mã đề.
    \item Tại bước 8, nếu phát hiện câu tô nhầm hoặc tô nhẹ: ứng dụng đánh dấu cảnh báo và yêu cầu giáo viên xác nhận.
    \item Tại bước 13, nếu mất kết nối mạng: ứng dụng lưu kết quả cục bộ và tự động đồng bộ khi có kết nối.
\end{itemize}

\noindent\textbf{Hậu điều kiện:} Nếu thành công, bài làm của học sinh được lưu trên server với điểm số đã tính, học sinh có thể xem điểm sau khi giáo viên phát hành kết quả. Nếu thất bại, ứng dụng giữ kết quả ở trạng thái chờ xử lý.

\subsection{Đặc tả use case Tạo đơn phúc khảo}
\label{subsection:2.3.4}

\noindent\textbf{Tên use case:} Tạo đơn phúc khảo.

\noindent\textbf{Tiền điều kiện:} Học sinh đã đăng nhập vào hệ thống, đã có kết quả bài thi được công bố và còn trong thời hạn phúc khảo.

\noindent\textbf{Luồng sự kiện chính:}
\begin{enumerate}
    \item Học sinh truy cập chức năng "Bài thi của tôi" và chọn bài thi muốn phúc khảo.
    \item Hệ thống hiển thị chi tiết bài thi gồm: điểm tổng, điểm từng câu, ảnh phiếu đã chấm và đáp án đúng.
    \item Học sinh chọn chức năng "Tạo đơn phúc khảo" và chọn một hoặc nhiều câu cần phúc khảo.
    \item Hệ thống hiển thị form yêu cầu nhập lý do phúc khảo cho từng câu, cho phép đính kèm hình ảnh minh chứng (nếu có).
    \item Học sinh nhập lý do, đính kèm bằng chứng (nếu có) và nhấn nút "Gửi đơn".
    \item Hệ thống kiểm tra thời hạn phúc khảo và tính hợp lệ của thông tin.
    \item Hệ thống tạo đơn phúc khảo mới với trạng thái "đang chờ", gửi thông báo cho giáo viên phụ trách.
    \item Hệ thống gửi thông báo xác nhận cho học sinh và hiển thị trạng thái đơn phúc khảo trong danh sách "Đơn phúc khảo của tôi".
\end{enumerate}

\noindent\textbf{Luồng sự kiện phát sinh:}
\begin{itemize}
    \item Tại bước 6, nếu hết thời hạn phúc khảo: hệ thống từ chối và thông báo thời hạn đã kết thúc.
    \item Tại bước 6, nếu học sinh đã gửi đơn cho cùng một câu: hệ thống từ chối và đề xuất xem trạng thái đơn cũ.
    \item Tại bước 4, nếu lý do phúc khảo quá ngắn: hệ thống yêu cầu nhập chi tiết hơn.
\end{itemize}

\noindent\textbf{Hậu điều kiện:} Nếu thành công, đơn phúc khảo được tạo và ở trạng thái "đang chờ", giáo viên nhận được thông báo để xử lý. Nếu thất bại, trạng thái các đơn cũ không thay đổi.

\subsection{Đặc tả use case Phản hồi phúc khảo}
\label{subsection:2.3.5}

\noindent\textbf{Tên use case:} Phản hồi phúc khảo.

\noindent\textbf{Tiền điều kiện:} Giáo viên đã đăng nhập, có đơn phúc khảo đang ở trạng thái "đang chờ".

\noindent\textbf{Luồng sự kiện chính:}
\begin{enumerate}
    \item Giáo viên truy cập chức năng "Phúc khảo" và xem danh sách đơn phúc khảo đang chờ.
    \item Giáo viên chọn một đơn phúc khảo để xem chi tiết.
    \item Hệ thống hiển thị thông tin bài thi của học sinh, câu được phúc khảo, lý do và hình ảnh minh chứng (nếu có).
    \item Giáo viên so sánh câu trả lời của học sinh với đáp án, đối chiếu với hình ảnh phiếu gốc.
    \item Giáo viên quyết định phê duyệt hoặc từ chối và nhập nội dung phản hồi cho học sinh.
    \item Giáo viên nhấn nút "Gửi phản hồi" để hoàn tất.
    \item Hệ thống cập nhật trạng thái đơn phúc khảo thành "đã phê duyệt" hoặc "đã từ chối".
    \item Nếu phê duyệt, hệ thống tính lại điểm cho câu được phúc khảo và cập nhật tổng điểm của học sinh.
    \item Hệ thống gửi thông báo kết quả phúc khảo cho học sinh.
\end{enumerate}

\noindent\textbf{Luồng sự kiện phát sinh:}
\begin{itemize}
    \item Tại bước 5, nếu giáo viên cần thêm thời gian: đơn được đánh dấu "đang xử lý" và có thể tiếp tục sau.
    \item Tại bước 8, nếu có lỗi khi tính lại điểm: hệ thống ghi nhận lỗi, giữ nguyên điểm cũ và thông báo cho quản trị hệ thống.
    \item Tại bước 5, nếu đơn đã hết thời hạn xử lý: hệ thống cảnh báo và yêu cầu giáo viên xác nhận trước khi gửi.
\end{itemize}

\noindent\textbf{Hậu điều kiện:} Nếu thành công, đơn phúc khảo được đóng và điểm được cập nhật (nếu phê duyệt). Học sinh nhận được thông báo kết quả.

\subsection{Đặc tả use case Tạo báo cáo cho bài kiểm tra}
\label{subsection:2.3.6}

\noindent\textbf{Tên use case:} Tạo báo cáo cho bài kiểm tra.

\noindent\textbf{Tiền điều kiện:} Giáo viên đã đăng nhập, bài kiểm tra đã có kết quả chấm từ học sinh.

\noindent\textbf{Luồng sự kiện chính:}
\begin{enumerate}
    \item Giáo viên truy cập chức năng "Báo cáo" và chọn bài kiểm tra cần tạo báo cáo.
    \item Hệ thống tổng hợp dữ liệu: danh sách bài làm, điểm số, đáp án, thời gian làm bài.
    \item Hệ thống (System) phân tích thống kê: tính phổ điểm, điểm trung bình, độ lệch chuẩn, tỉ lệ đạt, độ khó và độ phân biệt của từng câu hỏi.
    \item Hệ thống sinh báo cáo dưới dạng biểu đồ (cột, tròn, phân tán) và bảng số liệu.
    \item Hệ thống sử dụng AI để sinh nhận xét tổng quan về chất lượng đề thi và gợi ý cải thiện.
    \item Giáo viên xem trước báo cáo, có thể tùy chỉnh các chỉ số muốn hiển thị.
    \item Giáo viên chọn xuất báo cáo ra định dạng PDF hoặc Excel.
    \item Hệ thống sinh file báo cáo và cho phép giáo viên tải về.
\end{enumerate}

\noindent\textbf{Luồng sự kiện phát sinh:}
\begin{itemize}
    \item Tại bước 3, nếu dữ liệu chưa đầy đủ: hệ thống cảnh báo và chỉ sinh báo cáo với các chỉ số khả dụng.
    \item Tại bước 5, nếu dịch vụ AI gặp sự cố: hệ thống bỏ qua phần nhận xét AI và tiếp tục sinh phần còn lại của báo cáo.
    \item Tại bước 7, nếu xuất file thất bại: hệ thống thông báo lỗi và cho phép thử lại.
\end{itemize}

\noindent\textbf{Hậu điều kiện:} Nếu thành công, giáo viên nhận được file báo cáo hoàn chỉnh để sử dụng cho công tác đánh giá chất lượng giảng dạy và ra đề.

\section{Yêu cầu phi chức năng}
\label{section:2.4}

Ngoài các yêu cầu chức năng đã nêu, hệ thống Smart Grading cần đáp ứng các yêu cầu phi chức năng sau:

\noindent\textbf{Yêu cầu về hiệu năng.} Hệ thống phải đảm bảo thời gian phản hồi của các API dưới 2 giây đối với các thao tác thông thường và dưới 5 giây đối với các thao tác tổng hợp dữ liệu lớn như tạo báo cáo. Thời gian chấm một phiếu trả lời bằng ứng dụng di động không quá 3 giây với điều kiện ảnh đầu vào đạt chất lượng.

\noindent\textbf{Yêu cầu về độ tin cậy.} Hệ thống phải đảm bảo độ chính xác của thuật toán nhận dạng OMR đạt tối thiểu 95\% trên các điều kiện ảnh chụp thông thường. Hệ thống phải có cơ chế ghi log đầy đủ để truy vết lỗi và phải có khả năng khôi phục sau sự cố mà không mất dữ liệu đã chấm.

\noindent\textbf{Yêu cầu về tính dễ dùng.} Giao diện người dùng phải được Việt hóa hoàn toàn, sử dụng phông chữ dễ đọc và có độ tương phản phù hợp. Các thao tác chính phải được thực hiện trong vòng không quá 3 lần nhấp chuột hoặc chạm. Ứng dụng di động phải hoạt động ổn định trên các thiết bị có cấu hình trung bình với hệ điều hành Android từ 7.0 trở lên.

\noindent\textbf{Yêu cầu về tính bảo trì.} Mã nguồn phải được tổ chức theo kiến trúc phân lớp rõ ràng, có tài liệu hướng dẫn cài đặt và vận hành. Hệ thống phải có bộ kiểm thử tự động cho các thành phần cốt lõi với độ phủ kiểm thử tối thiểu 60\%.

\noindent\textbf{Yêu cầu về bảo mật.} Hệ thống phải sử dụng JWT cho xác thực, mật khẩu phải được lưu trữ dưới dạng mã băm bcrypt. Các điểm cuối API phải được phân quyền theo vai trò người dùng. Ảnh phiếu thi và dữ liệu cá nhân của học sinh phải được bảo vệ khỏi truy cập trái phép.

\noindent\textbf{Yêu cầu về cơ sở dữ liệu và công nghệ.} Hệ thống sử dụng cơ sở dữ liệu MongoDB phiên bản 5.0 trở lên để lưu trữ dữ liệu nghiệp vụ. Hệ thống phải hỗ trợ triển khai trên cả môi trường phát triển (local) và môi trường đám mây (cloud). Backend chạy trên Node.js v18 trở lên, frontend web xây dựng trên React 18 với TypeScript, ứng dụng di động xây dựng trên Flutter 3.x.

\section{Kết chương}
Chương 2 đã trình bày kết quả khảo sát hiện trạng từ ba nguồn: nhu cầu người dùng, hệ thống chấm thi hiện tại tại các trường và các ứng dụng tương tự trên thị trường. Trên cơ sở đó, chương này đã mô tả tổng quan chức năng của hệ thống Smart Grading với mười ba use case chính phân thành bốn nhóm, đồng thời đặc tả chi tiết sáu use case quan trọng nhất với đầy đủ luồng sự kiện chính, luồng phát sinh, tiền điều kiện và hậu điều kiện. Bên cạnh đó, các yêu cầu phi chức năng về hiệu năng, độ tin cậy, tính dễ dùng, bảo trì, bảo mật và công nghệ cũng đã được xác định rõ. Các kết quả chính của chương gồm:
\begin{itemize}
    \item bảng so sánh tính năng giữa các giải pháp hiện có;
    \item biểu đồ use case tổng quát và mười hai biểu đồ phân rã cho từng nhóm use case;
    \item ba quy trình nghiệp vụ trọng yếu;
    \item sáu đặc tả use case chi tiết.
\end{itemize}
Chương 3 sẽ trình bày các nền tảng lý thuyết và công nghệ được lựa chọn để hiện thực hóa các yêu cầu đã nêu trong chương này.

\end{document}